\documentclass[12pt]{article}
\usepackage[utf8]{inputenc}
\usepackage[T1]{fontenc}
\usepackage{hyperref}
\title{Meriam Mir Language Resources for K-12 Curriculum Development}
\author{Compiled by Claude Sonnet 5 (claude-sonnet-5) for Andr\'as Kornai}
\date{}
\begin{document}
\maketitle

\section{Introduction}
Meriam Mir (ISO 639-3 ulk) is an endangered Eastern Trans-Fly language of the eastern Torres Strait, Australia. The 2016 Australian census reportedly recorded 217 speakers (as summarised in one search result). Searches found no school curriculum or online course in the language. What exists is mostly dictionaries, word lists and archival manuscripts.

\section{Primary Educational Resources}
No Meriam Mir primer, textbook or classroom programme turned up in the searches. The nearest items are the following.
\begin{itemize}
\item State Library of Queensland, Torres Strait Everyday Words (word list covering both Torres Strait languages and Torres Strait Creole; the library notes that spellings and pronunciations vary and advises checking with local speakers): \url{https://www.slq.qld.gov.au/sites/default/files/Torres%20Strait%20Everyday%20Words.pdf}
\item State Library of Queensland, Languages of the Torres Straits (PDF): \url{https://content.slq.qld.gov.au/sites/default/files/Languages%20of%20the%20Torres%20Straits_0.pdf}
\item State Library of Queensland blog, Word Week (2019 International Year of Indigenous Languages): \url{https://www.slq.qld.gov.au/blog/2019-international-year-indigenous-languages-word-week-week-thirty-three}
\end{itemize}

\section{Linguistic and Reference Materials}
\begin{itemize}
\item Anna Shnukal (1992), Meriam Mir lexicon in Torres Strait Creole (Glottolog bibliographic record): \url{https://glottolog.org/resource/reference/id/305459}
\item Sydney H. Ray (2003), Dictionary of Torres Strait languages, 2nd edition (Glottolog bibliographic record): \url{https://glottolog.org/resource/reference/id/6723}
\item Glottolog entry for Meriam Mir: \url{https://glottolog.org/resource/languoid/id/meri1244}
\item Wikipedia, Meriam Mir: \url{https://www.Wikipedia.org/wiki/Meriam_Mir}
\end{itemize}
The search summary also mentioned McConvell, Day and Black (1983), ``Making a Meriam Mir Dictionary'', and Babia and Day (1989), Torres Strait Picture Dictionary. No URL for either appeared in the results, so they are not linked.

\section{Religious and Community Resources}
No scripture translations or community organisation pages specific to Meriam Mir were found. The Torres Strait Traditional Languages Centre was named in a search summary as the regional body for language work, but no URL for it was seen, so it is not linked.

\section{Digital and Audio Resources}
\begin{itemize}
\item University of Queensland Library, Torres Strait Islanders -- Languages (handwritten manuscripts, including a Mer vocabulary and a ``Mer -- Grammar'' notebook by Robert Victor Miles): \url{https://manuscripts.library.uq.edu.au/index.php/torres-strait-islanders-languages}
\item Endangered Languages Project entry: \url{https://www.endangeredlanguages.com/bn/node/95993}
\end{itemize}
No audio, video or app resources for Meriam Mir were found.

\section{Curriculum Development Context}
Little is available in print or online. Any curriculum would have to start from the dictionaries and archival manuscripts above and from work with the remaining speakers. The community and the regional language centre should lead that work, and the word lists' own advice on local pronunciation should be followed. The dossier does not claim any existing school programme.

\end{document}